%! The Fundamental Theorem of Linear Algebra — Unit 3, Lesson 3.6 — Lesson Plan
\documentclass[10pt]{article}
\usepackage{linalg-article}
\usepackage{linalg-boxes}
\usepackage{graphicx}   % -article does NOT load graphicx; needed for thumbnails/screenshots
\graphicspath{{images/}}
\newcommand{\TallMath}[1]{$\displaystyle #1\rule[-1.4em]{0pt}{3.2em}$}
\newcommand{\vv}{\mathbf{v}}
\newcommand{\xx}{\mathbf{x}}
\newcommand{\yy}{\mathbf{y}}
\newcommand{\bb}{\mathbf{b}}
\newcommand{\svec}{\mathbf{s}}
\newcommand{\zero}{\mathbf{0}}
\newcommand{\T}{\mathsf{T}}
\newcommand{\UnitNumberName}{Unit 3: The Four Fundamental Subspaces \quad}
\newcommand{\LessonNumberName}{Lesson 3.6: The Fundamental Theorem of Linear Algebra}

\begin{document}
\begin{center}
    {\Large\bfseries \CourseName: \SchoolYear \\
    {\small\bfseries \UnitNumberName \LessonNumberName}}
\end{center}

\begin{tcolorbox}[colback=blush, colframe=burgundy, boxrule=0.9pt, arc=2mm,
    left=2mm, right=2mm, top=1.5mm, bottom=1.5mm]
    \textbf{Primary Objective:} Students can assemble the four fundamental subspaces of an $m\times n$ matrix
    $A$ of rank $r$ into the \emph{big picture of linear algebra}: the input room $\mathbb{R}^n$ split into
    the row space $C(A^{\T})$ (dimension $r$) and nullspace $N(A)$ (dimension $n-r$), the output room
    $\mathbb{R}^m$ split into the column space $C(A)$ (dimension $r$) and left nullspace $N(A^{\T})$
    (dimension $m-r$). They can state \textbf{Part~1} (those dimensions and the counting rules) and
    \textbf{Part~2} (the two subspaces in each room are \emph{orthogonal complements}: the row space is
    perpendicular to the nullspace, the column space to the left nullspace), and justify Part~2 from a single
    idea --- $A\xx=\zero$ means every row of $A$ dotted with $\xx$ is $0$, so $\xx$ is perpendicular to the
    whole row space.
\end{tcolorbox}

\renewcommand{\arraystretch}{1.4}
\begin{skillbox}[Priority Ideas \& Skills]{goldbox}
\begin{tabularx}{\linewidth}{@{}X|X@{}}
\textbf{Priority Ideas \& Skills} & \textbf{Key Understandings (the ``why'')} \\[2pt]
\hline
\vspace{-6pt}
\begin{itemize}[leftmargin=1.1em, nosep, after=\vspace{2pt}]
  \item Draw and label the \textbf{big-picture diagram}: two rooms, four subspaces, four dimensions.
  \item State \textbf{Part~1} (dimensions $r,\,n-r,\,r,\,m-r$) and \textbf{Part~2} (orthogonality) of the theorem.
  \item Show $A\xx=\zero \Rightarrow \xx\perp$ every row, so $N(A)\perp C(A^{\T})$; likewise $N(A^{\T})\perp C(A)$.
  \item Verify a pair is orthogonal by checking \textbf{dot products $=0$} on a basis.
\end{itemize}
&
\vspace{-6pt}
\begin{itemize}[leftmargin=1.1em, nosep, after=\vspace{2pt}]
  \item $A\xx$ is computed \emph{row by row}, so $A\xx=\zero$ literally says each row $\cdot\,\xx=0$.
  \item ``Perpendicular $\Leftrightarrow$ dot product $0$'' (Lesson~1.2) is the bridge from subspaces to right angles.
  \item In each room the two subspaces meet only at $\zero$ and, being perpendicular, \emph{fill} the room ($r+(n-r)=n$).
  \item The theorem is the whole unit in one picture --- and the door into orthogonality (Unit~4).
\end{itemize}
\end{tabularx}
\end{skillbox}

\begin{skillbox}[{Vocabulary, Concepts, \& Theorems}]{greenbox}
\renewcommand{\arraystretch}{1.3}
\begin{tabularx}{\linewidth}{@{}l X@{}}
\textbf{Fundamental Theorem (Part 1)} & Dimensions: $\dim C(A)=\dim C(A^{\T})=r$, $\dim N(A)=n-r$, $\dim N(A^{\T})=m-r$. \\
\textbf{Fundamental Theorem (Part 2)} & Orthogonality: $N(A)\perp C(A^{\T})$ in $\mathbb{R}^n$; $N(A^{\T})\perp C(A)$ in $\mathbb{R}^m$. \\
\textbf{Orthogonal} & Two vectors are orthogonal (perpendicular) when their dot product is $0$ (Lesson~1.2). \\
\textbf{Orthogonal complement} & Two subspaces that are perpendicular \emph{and} whose dimensions add to fill the room --- each is ``everything perpendicular to the other.'' \\
\textbf{Big picture} & The two-room diagram holding all four subspaces, their dimensions, and the right angles. \\
\end{tabularx}
\end{skillbox}

\begin{fixedskillbox}[Activate Prior Knowledge \& Spiral Review]{blush}
    The Warm-Up rehearses the two tools this capstone fuses: \textbf{(1)} the \emph{dot-product test for
    perpendicularity} (Lesson~1.2) --- two vectors are perpendicular exactly when their dot product is $0$;
    and \textbf{(2)} reading $A\xx$ \emph{row by row} (Lesson~1.3) --- each entry of $A\xx$ is a row of $A$
    dotted with $\xx$. Debrief the punchline aloud: putting those together, $A\xx=\zero$ says every row of
    $A$ is perpendicular to $\xx$ --- which is exactly why the nullspace sits at right angles to the row
    space. Item~3 recalls the four dimensions from Lesson~3.5, the ``Part~1'' the lesson assembles.
\end{fixedskillbox}

\begin{skillbox}[Hook]{blush}
    One last visit to the robot arm --- to see the whole machine at once. Across the unit we met the four
    subspaces one at a time: the reachable tip positions (column space), the do-nothing lever settings
    (nullspace), and their transposes (row space, left nullspace). \textbf{``What does the whole matrix look
    like in one picture?''} Draw the \emph{two rooms} --- input $\mathbb{R}^n$ and output $\mathbb{R}^m$ ---
    and drop the four subspaces into place: two per room, dimensions $r$, $n-r$, $r$, $m-r$. Then the
    surprise: within each room the two subspaces are not just complementary in size --- they sit at a
    \textbf{right angle}. The do-nothing settings are perpendicular to the settings that ``really move'' the
    tip. That picture, dimensions plus right angles, \emph{is} the Fundamental Theorem of Linear Algebra.
\end{skillbox}

\begin{skillbox}[Lesson]{blush}
\begin{multicols}{2}
\textbf{1. Assemble the big picture (Part 1).} An $m\times n$ matrix connects two rooms: input $\mathbb{R}^n$
(holding the row space $C(A^{\T})$ and nullspace $N(A)$) and output $\mathbb{R}^m$ (holding the column space
$C(A)$ and left nullspace $N(A^{\T})$). From one number $r$ come all four dimensions: $r$, $n-r$, $r$, $m-r$,
with $r+(n-r)=n$ and $r+(m-r)=m$. This is Lesson~3.5, now drawn as \emph{one} diagram.

\textbf{2. The right-angle surprise (Part 2).} Compute $A\xx$ row by row: entry $i$ is (row$_i$)$\cdot\xx$.
So $A\xx=\zero$ means \emph{every} row of $A$ is perpendicular to $\xx$ --- hence $\xx$ is perpendicular to
every combination of rows, the entire row space. Therefore $N(A)\perp C(A^{\T})$.

\columnbreak

\textbf{3. The other pair, by transpose.} Apply the same idea to $A^{\T}$: $A^{\T}\yy=\zero$ means $\yy$ is
perpendicular to every row of $A^{\T}$ --- i.e.\ every \emph{column} of $A$. So $N(A^{\T})\perp C(A)$. Each
room holds a perpendicular pair.

\textbf{4. Orthogonal complements.} In $\mathbb{R}^n$ the row space and nullspace are perpendicular
\emph{and} their dimensions add to $n$ --- they meet only at $\zero$ and together fill the room. Each is
``everything perpendicular to the other'': \emph{orthogonal complements}. Same story in $\mathbb{R}^m$ for
column space and left nullspace. Part~1 (dimensions) $+$ Part~2 (right angles) $=$ the Fundamental Theorem,
and the launch pad for orthogonality in Unit~4.
\end{multicols}
\end{skillbox}

\begin{skillbox}[Explicit Instruction: Why the Nullspace $\perp$ the Row Space]{blush}
\begin{tabularx}{\linewidth}{@{}p{0.46\linewidth} X@{}}
\textbf{The one-line argument:}
\begin{enumerate}[leftmargin=1.3em, nosep]
  \item $A\xx$ is a stack of dot products: entry $i$ is (row$_i$ of $A$)$\,\cdot\,\xx$.
  \item $A\xx=\zero$ $\Rightarrow$ every row$_i\cdot\xx=0$.
  \item Dot $=0$ means perpendicular (Lesson~1.2), so $\xx\perp$ each row.
  \item Perpendicular to each row $\Rightarrow$ perpendicular to every combination of rows $=$ the row space.
  \item So $N(A)\perp C(A^{\T})$. (Transpose for $N(A^{\T})\perp C(A)$.)
\end{enumerate}
&
\vspace{-2pt}
\textbf{Worked} ($A=\begin{bsmallmatrix} 1&1&2\\2&1&3\\3&2&5 \end{bsmallmatrix}$, $r=2$). Nullspace vector
$\svec=\begin{bsmallmatrix} -1\\-1\\1 \end{bsmallmatrix}$; row space basis $(1,0,1),(0,1,1)$.
Check: $(1,0,1)\cdot\svec=-1+0+1=0$ and $(0,1,1)\cdot\svec=0-1+1=0$ --- $\svec$ is perpendicular to the whole
row space. Left-null $\yy=\begin{bsmallmatrix} 1\\1\\-1 \end{bsmallmatrix}$; columns
$(1,2,3),(1,1,2)$: $\yy\cdot(1,2,3)=0$, $\yy\cdot(1,1,2)=0$ --- $\yy\perp C(A)$.
\end{tabularx}
\end{skillbox}

\begin{skillbox}[Active Monitoring]{redbox}
Circulate during the notes practice and Tier work. Watch for and correct:
\begin{itemize}[leftmargin=1.2em, nosep]
  \item pairing the wrong subspaces --- the right angle is \emph{within} a room (row space~$\perp$~nullspace in $\mathbb{R}^n$; column space~$\perp$~left nullspace in $\mathbb{R}^m$), never across rooms;
  \item testing orthogonality against only \emph{one} row/column --- perpendicular to the space means dot $=0$ with \emph{every} basis vector;
  \item confusing ``dimensions add to $n$'' (Part~1) with ``perpendicular'' (Part~2) --- orthogonal complements need \emph{both};
  \item forgetting the argument reduces to Lesson~1.2: dot product $0$ \emph{is} perpendicularity --- no new machinery.
\end{itemize}
Cold-call prompts: ``When you compute $A\xx$, what is entry $i$?'' ``So what does $A\xx=\zero$ say about the
rows and $\xx$?'' ``Which two subspaces are perpendicular in $\mathbb{R}^m$?'' ``What two things does
`orthogonal complement' require?''
\end{skillbox}

\begin{skillbox}[Group Work \& Differentiation]{redbox}
\textbf{The Big Picture} activity (tiered; mirrors the notes).
\begin{multicols}{3}
\textbf{Tier R --- Remediate.} For a matrix with $m$, $n$, $r$ given: fill in the big-picture diagram ---
place each subspace in its room and label its dimension; name the two perpendicular pairs.

\columnbreak
\textbf{Tier A --- Approaching.} For a concrete matrix (reduction supplied): list a basis for all four
subspaces, then \emph{verify} both right angles by computing dot products against the bases.

\columnbreak
\textbf{Tier E --- Extension.} Explain why $A\xx=\zero$ forces $\xx\perp$ the row space; why the two
perpendicular subspaces meet only at $\zero$; and why their dimensions must add to fill the room.
\end{multicols}
\end{skillbox}

\begin{skillbox}[Individual Work \& Assessment]{redbox}
\textbf{Exit Ticket} (independent, no notes): fill in a big-picture diagram (rooms, subspaces, four
dimensions); verify one right angle by a dot-product check; and a \textbf{conceptual/justification check} ---
explain, from $A\xx=\zero$, why the nullspace is perpendicular to the row space. Collect and sort into ``got
it / paired across rooms / checked only one vector'' to gauge readiness for the Unit~3 test and Unit~4.
\end{skillbox}

\begin{skillbox}[Reinforcement \& Extension]{goldbox}
\textbf{Homework:} state both parts of the theorem; complete big-picture diagrams; verify orthogonality by
dot products; and read the right angles back as statements about rows and columns. \textbf{Extension:} why
orthogonal complements meet only at $\zero$; the invertible case (both nullspaces $\{\zero\}$, both spaces
fill $\mathbb{R}^n$). \textbf{Preview:} the Unit~3 test, then \emph{Unit~4: Orthogonality} --- projecting
onto subspaces and least squares, built directly on the right angles discovered here.
\end{skillbox}
\end{document}
